\begin{titlepage}
\thispagestyle{empty}
\begin{tikzpicture}[remember picture,overlay]
	% Grundfläche
	\fill[SmartPaper] (current page.south west) rectangle (current page.north east);

	% Diagonale Trennung: helle obere linke Fläche
	\fill[white] (current page.north west)
		-- ($(current page.north east)+(0,-8.8cm)$)
		-- ($(current page.south east)+(-8.8cm,0)$)
		-- (current page.south west)
		-- cycle;

	% Dunklere technische Fläche unten rechts
	\fill[SmartBlue!92] (current page.south east)
		-- ($(current page.south west)+(0,8.8cm)$)
		-- ($(current page.north west)+(8.8cm,0)$)
		-- (current page.north east)
		-- cycle;

	% Feine Diagonallinie
	\draw[SmartTeal!55,line width=1.2pt] ($(current page.north east)+(0,-8.8cm)$) -- ($(current page.south west)+(8.8cm,0)$);

	% Titelblock mit transparenter Fläche für hohe Lesbarkeit.
	\node[anchor=north west, rounded corners=3mm, fill=white, fill opacity=0.86, text opacity=1,
		inner xsep=18pt, inner ysep=16pt, text width=0.42\paperwidth, align=left]
		at ($(current page.north west)+(1.35cm,-1.55cm)$) {%
		\begin{minipage}[t]{0.42\paperwidth}
			{\sffamily\bfseries\small\color{SmartTeal}TECHNISCHES HANDBUCH}\\[1.35em]
			{\sffamily\bfseries\color{SmartInk}\fontsize{31}{35}\selectfont Handbuch zum\\(be)greifbaren Smart Home}\\[1.0em]
			{\sffamily\normalsize\color{SmartInk}Konzeption und Konstruktion einer modularen 12V-IoT-Lernumgebung für die Digitale Grundbildung}\\[1.45em]
			\CoverLogoGraphic\\[1.55em]
			{\sffamily\normalsize\color{SmartInk}Sekundarstufe I \textbullet{} Unterrichts- und Laborumgebung \textbullet{} modulare Hardware \textbullet{} Digitale Grundbildung}
		\end{minipage}
	};

	% Bildbereich mit sehr leichter Fläche, damit das Platzhalterbild ruhig wirkt.
	\node[anchor=north east, rounded corners=3mm, fill=white, fill opacity=0.14, text opacity=1,
		inner xsep=14pt, inner ysep=14pt, text width=0.25\paperwidth, align=center]
		at ($(current page.north east)+(-1.45cm,-2.00cm)$) {
			\CoverImageGraphic
		};

	% Unterer Abschlussbereich für Autor und Erscheinungsjahr.
	\node[anchor=south, rounded corners=3mm, fill=SmartLight, fill opacity=0.92, text opacity=1,
		inner xsep=18pt, inner ysep=12pt, text width=0.82\paperwidth, align=center]
		at ($(current page.south)+(0,1.15cm)$) {%
		\begin{minipage}[t]{0.82\paperwidth}
			\begin{tabular*}{\linewidth}{@{}l@{\extracolsep{\fill}}r@{}}
				{\sffamily\bfseries\color{SmartBlue}\small Autor} & {\sffamily\bfseries\color{SmartTeal}\small ERSCHEINUNGSJAHR}\\[0.65em]
				{\sffamily\bfseries\color{SmartInk}\fontsize{20}{22}\selectfont Johannes Baischer} & {\sffamily\bfseries\color{SmartInk}\fontsize{20}{22}\selectfont 2026}
			\end{tabular*}
		\end{minipage}
		};
\end{tikzpicture}
\end{titlepage}